% Generated draft: requires mathematical and source review.
\documentclass[11pt]{article}
\usepackage[margin=25mm]{geometry}
\usepackage{fontspec}
\setmainfont{Latin Modern Roman}
\usepackage{amsmath,amssymb,mathtools,unicode-math}
\setmathfont{Latin Modern Math}
\usepackage{enumitem,needspace,xurl,hyperref}
\setlength{\parindent}{0pt}
\setlength{\parskip}{5pt}
\setlength{\emergencystretch}{3em}
\newcommand{\N}{\mathbb{N}}
\newcommand{\Z}{\mathbb{Z}}
\newcommand{\Q}{\mathbb{Q}}
\newcommand{\R}{\mathbb{R}}
\newcommand{\C}{\mathbb{C}}
\newcommand{\F}{\mathbb{F}}
\newcommand{\A}{\mathbb{A}}
\newcommand{\PP}{\mathbb{P}}
\newcommand{\E}{\mathbb{E}}
\newcommand{\eps}{\varepsilon}
\newcommand{\dd}{\,\mathrm{d}}
\begin{document}
\section*{UnsolvedMath ID: 30003976}
\textbf{Source problem code:} catalogue entry
\subsection{Definitions and mathematical statement}
Let \(V\) be a finite-dimensional complex vector space. A complex reflection is a nonidentity finite-order linear transformation fixing a hyperplane pointwise. A finite complex reflection group \(G\subseteq\operatorname{GL}(V)\) is a finite group generated by complex reflections; it is irreducible if \(V\) has no nonzero proper \(G\)-invariant linear subspace. Its reflection arrangement is
\[
\mathcal A(G)=\{\ker(s-\operatorname{id}_V):s\in G\text{ is a complex reflection}\},
\]
where repeated hyperplanes are counted only once.

The supplied question asks for a classification of the irreducible complex reflection groups whose triple-point ideal satisfies
\[
I_X^{(3)}\not\subseteq I_X^2.
\]
However, the source record does not define the triple-point locus \(X\), specify the dimension of \(V\), or specify whether triple points means points on exactly three or on at least three arrangement components. These qualifications remain unresolved.

For a precise conditional formulation in dimension \(3\), projectivize the arrangement to a set of lines in \(\PP(V)\cong\PP^2_{\C}\), and put
\[
m_G(p)=\#\{H\in\mathcal A(G):p\in\PP(H)\}.
\]
Two possible reduced point loci are
\[
X_{=3}(G)=\{p\in\PP(V):m_G(p)=3\},\qquad
X_{\geq3}(G)=\{p\in\PP(V):m_G(p)\geq3\}.
\]
The record does not select between them. For either specified choice \(X\), let \(S=\operatorname{Sym}(V^*)\cong\C[x_0,x_1,x_2]\), and let \(\mathfrak p_p\subset S\) be the homogeneous prime ideal of the point \(p\). Under the reduced-point interpretation, define
\[
I_X=\bigcap_{p\in X}\mathfrak p_p,\qquad
I_X^{(3)}=\bigcap_{p\in X}\mathfrak p_p^3,\qquad
I_X^2=(fg:f,g\in I_X).
\]
An empty intersection is interpreted as \(S\). For each fixed choice of \(X\), the conditional classification problem is to determine all irreducible finite complex reflection groups \(G\subseteq\operatorname{GL}_3(\C)\) for which there exists
\[
f\in\bigcap_{p\in X(G)}\mathfrak p_p^3\quad\text{with}\quad f\notin I_{X(G)}^2.
\]
This dimension-three formulation is conditional: the supplied record does not authorize restricting its requested classification of all irreducible complex reflection groups to dimension three. A definition of the intended triple-point locus in other dimensions, and confirmation of its scheme structure, are needed to give the full source question a unique precise formulation.
\subsection{Short English statement}
The question asks which irreducible complex reflection groups give a reflection arrangement whose triple-point ideal has third symbolic power not contained in its ordinary square. The supplied record leaves the dimension and meaning of triple points unspecified. The plane-arrangement formulation above is therefore conditional and requires human mathematical review.
\subsection{Sources}
\begin{itemize}
\item Supplied local UnsolvedMath registry record, permanent ID \(30003976\), source problem code \texttt{catalogue entry}, titled \emph{Symbolic-Power Noncontainments for Reflection Arrangements}, in the Oberwolfach Reports Open Problems collection. No individual proposer is supplied; the proposed year is recorded as 2018.
\item Supplied report citation: \emph{Mini-Workshop: Asymptotic Invariants of Homogeneous Ideals (2019)}, Oberwolfach Reports, pp.\ 2739--2767, DOI \texttt{10.4171/owr/2018/45}. OWR report: \url{https://doi.org/10.4171/owr/2018/45}. The supplied citation labels the workshop 2019, while the DOI and proposed-year metadata indicate 2018; this discrepancy is unresolved.
\item Supplied literature source, labeled \emph{Reflection arrangements}: \url{https://arxiv.org/abs/2002.05353}.
\item Supplied literature source, labeled \emph{OWR 2022 report}: \url{https://ems.press/content/serial-article-files/46948}.
\item Literature assessment checked at \texttt{2026-08-21}: the supplied curation labels the problem open and reports that several reflection-group families yield noncontainments, while the classification of all irreducible complex reflection groups remains open. This is supplied dated curation derived from the OWR material, not an independently verified assessment. The supplied background explicitly requests independent verification of time-sensitive status and cited claims.
\item The registry labels the statement \texttt{corrected\_verified}, explaining that mathematical content was preserved while notation, typography, and sentence structure were repaired. This supplied editorial label does not establish mathematical verification of this draft.
\end{itemize}
\end{document}
